\section{Related work and software}
\label{sec:related}

\paragraph{Programs for discrete partial identification.}
\citet{balke1994bounds,balke1997bounds} wrote the instrumental-variable
problem as an LP over response functions and solved it symbolically.
Canonical SCMs extend the response-function parameterization to any
diagram \citep{zhang2022partial}, at the price of a polynomial program in
general; \citet{zhang2022partial} approach it by Bayesian sampling, and
\citet{duarte2024automated} by polynomial programming that returns sharp
bounds or valid outer bounds with a certified gap.
\citet{finkelstein2020bounds} derive bounds from logical relations among
counterfactuals. These works either fix one tractable design or give a sufficient condition
for a linear program. \Cref{thm:lp} is necessary and sufficient for diagrams
in which every $c$-component has at most two hyperedges: it says when the
program is an LP and when it is not, and the package uses that answer to pick
the solver and to label the result.

The same canonical model is also used without forming the program.
\citet{bellot2026foundation} put a full-support prior on canonical SCMs with discrete observables and train a prior-data fitted network \citep{muller2022transformers} that maps an observational table to a distribution over a counterfactual query.
The support of that distribution approximates the identified set.
\Cref{thm:lp} instead reads off the diagram and the query when the set is the optimum of one linear program.
Partial identification is also posed across domains.
\citet{jalaldoust2024partial} bound a functional of a target law, such as the risk of a fixed classifier, from source-domain data and a selection diagram, by gradient optimization over a structural parameterization.
The constraints there are invariances of mechanisms across domains, rather than the single observational table of \cref{def:reduced}.

\paragraph{Assumptions and sensitivity models.}
Shape restrictions \citep{manski1990nonparametric,manski1997monotone,manski2000monotone}
and sensitivity models \citep{rosenbaum2002observational,tan2006distributional,masten2018identification,jin2023sensitivity}
usually come with their own derivation and estimator, often for one
graph; see \citet{tamer2010partial} and \citet{kline2023recent} for reviews.
In the program of \cref{def:optimal_bound}, the anchored versions of these
models are linear constraints on the same response-type law
(\cref{prop:assumptions}), which is what lets the package compose them with each other
and with a diagram. Unanchored odds-ratio models and information-theoretic
budgets \citep{jung2026information} are not linear in this sense and are
outside the LP path.

\begin{table}[t]
\centering
\caption{Software for bounds on causal effects with discrete data.}
\label{tab:software}
\small
\begin{tabular}{@{}*{4}{>{\raggedright\arraybackslash}p{0.225\textwidth}}@{}}
\toprule
Tool & Language & Scope & Output \\
\midrule
\texttt{causaloptim} \newline \citep{symbolicbounds2023general} & R
  & graphs in their class, linear constraints & symbolic sharp bounds \\
Autobounds \newline \citep{duarte2024automated} & Python
  & any discrete graph, polynomial constraints & numeric bounds, optimality gap \\
CausalQueries \newline \citep{humphreys2023integrated} & R
  & discrete graphs, type restrictions & posterior over types \\
\texttt{bpbounds} \newline \citep{palmer2011nonparametric} & R, Stata
  & instrumental variables & closed-form Balke--Pearl bounds \\
\texttt{sensemakr} \newline \citep{cinelli2020making} & R, Stata, Python
  & linear regression & omitted-variable sensitivity \\
\pid{} & Python
  & discrete graphs or orders, named assumptions & numeric bounds, sharpness flag \\
\bottomrule
\end{tabular}
\end{table}

\paragraph{Software.}
\Cref{tab:software} separates tools that bound a counterfactual query in a
discrete structural model from tools that answer a different question.
\texttt{bpbounds} returns the closed-form Balke--Pearl interval for a binary
instrument. CausalQueries returns a posterior over response types.
\texttt{sensemakr} reports how a coefficient in a linear regression moves
when an omitted variable of a stated strength is introduced. The other three
packages take a discrete diagram and a causal query and return bounds.

\texttt{causaloptim} enumerates the vertices of one linear program and
returns a symbolic expression in the observational probabilities. The
expression applies to every table, for graphs that meet their
Proposition~2 and for constraints that stay linear in one joint type law.
\Cref{thm:lp} contains that class (\cref{sec:sachs}). Inside diagrams with
at most two hyperedges in each component, \cref{thm:lp} is an equivalence
and their Proposition~2 is a sufficient condition. \Cref{app:sachs} lists the
graphs on which \cref{thm:lp} is a linear program and their Proposition~2 is not.

Autobounds accepts an arbitrary discrete diagram and encodes the canonical
model as a polynomial program. \Cref{thm:lp} decides, before a solver is
called, when that program is linear. \Cref{app:autobounds} runs both packages
on the examples of \cref{sec:illustrations}.

\subsection{Comparison with Sachs et al.}
\label{app:sachs}

\citet[Proposition~2]{symbolicbounds2023general} give a sufficient condition
for a tight linear program in one joint type law, and \texttt{causaloptim}
returns a symbolic bound for the graphs that meet it. \Cref{sec:sachs}
records the four conditions and shows that every graph meeting them meets
\cref{thm:lp}. Their condition is sufficient. \Cref{thm:lp} is an equivalence
for diagrams whose components have at most two hyperedges, and the inclusion
of their class is one way. The rows below are the other direction: \cref{thm:lp} returns a linear program, and no
bipartition meets their conditions. The comparison is the criterion, not the
running time. The classical instrument is the reference case, in both
classes.

\begin{table}[ht]
\centering
\caption{Graphs on which \cref{thm:lp} is a linear program.
``In both'' means Proposition~2 of \citet{symbolicbounds2023general} applies
as well. The other rows have no bipartition meeting their conditions
(1)--(4). Intervals are the sharp bounds under \cref{thm:lp}.}
\label{tab:sachs}
\small
\begin{tabular}{@{}p{0.28\textwidth}lp{0.32\textwidth}l@{}}
\toprule
Graph & Prop.~2 & Why it fails & Interval \\
\midrule
Instrument, $X\leftrightarrow Y$ & in both & --- & $[0.283,0.615]$ \\
Mediated $X\leftrightarrow W\to Y$ & no & $Y$ is pinned, not confounded with $\{X,W\}$ & $[0.470,0.570]$ \\
Two instruments $Z_1,Z_2$ & no & $Z_1$ and $Z_2$ share no latent & $[0.397,0.789]$ \\
Instrument with $Z\leftrightarrow X$ & no & the latent on $\{Z,X\}$ crosses every cut & $[0.283,0.615]$ \\
Pair chain, PNS & no & two latents in one $c$-component & $[0.000,0.600]$ \\
\bottomrule
\end{tabular}
\end{table}

\paragraph{Mediated chain.}
The graph is $X\to W\to Y$ with one latent on $\{X,W\}$ and $U\perp U_Y$,
the query is $P(Y_1{=}1)$, and the table is~\eqref{eq:re3}. The event uses
$Y$ in one world, so \cref{lem:middle} pins $Y$ and \cref{thm:lp} leaves an
LP in $P(U)$ (\cref{ex:mediated}). The three bipartitions of $\{X,W,Y\}$ each
fail their conditions: $Y$ does not share the latent of $\{X,W\}$, the latent
on $\{X,W\}$ meets both sides of the cut after $X$, and $P(Y_x)$ is not a
linear functional of the types of $Y$ alone. Their algorithm therefore does
not apply. Forcing it by taking a free joint over $(R_X,R_W,R_Y)$ drops
$U\perp U_Y$ and widens the interval from $[0.470,0.570]$ to
$[0.270,0.770]$.

\paragraph{Two instruments.}
$Z_1,Z_2\to X\to Y$ with $X\leftrightarrow Y$ and $Z_1\perp Z_2$. The only
cut with all arrows pointing the same way puts both instruments on the left
and $\{X,Y\}$ on the right. Their condition~(3) asks the left side to share
one latent, and $Z_1$ and $Z_2$ do not. Both instruments are $c$-components
the query does not read, so \cref{thm:lp} still reduces the problem to the
Balke--Pearl LP in the response types of $\{X,Y\}$. On the packaged table the
interval is $[0.397,0.789]$.

\paragraph{A latent between the instrument and the treatment.}
Add the hyperedge $\{Z,X\}$ to the classical instrument, and keep
$\{X,Y\}$. The query $P(Y_1{=}1)$ does not change. Every bipartition either
lets the latent on $\{Z,X\}$ or the latent on $\{X,Y\}$ meet both sides, or
asks $\{Z,X,Y\}$ to share one latent. Proposition~2 does not apply.
\Cref{prop:lift} does: $Z$ is a root, the independence of the two latents is
linear, and the sharp interval equals the Balke--Pearl interval of the
classical instrument. On the seed-$2$ table both graphs give
$[0.283,0.615]$. Sharing a latent with the treatment does not change the
bound, and it takes the graph out of their class.

\paragraph{Pair chain.}
$X\to W\to Y$ with hyperedges $\{X,W\}$ and $\{W,Y\}$ is one $c$-component
with two latents (\cref{ex:pair}). The same case analysis leaves no
bipartition: a single latent does not cover all three nodes, and each
pairwise latent crosses the cut that would separate its endpoints.
$X$ is a root, so \cref{lem:two} makes the sharp program an LP for every
linear query, including PNS. On the table~\eqref{eq:re2} that LP gives
$[0.000,0.600]$. One latent on $\{X,W,Y\}$, which does meet their
conditions, gives the wider interval $[0.000,0.650]$. For $P(Y_1{=}1)$ on
this table the two models agree, at $[0.350,0.850]$; the probabilities of
causation are where the second latent binds.

\subsection{Comparison with Autobounds}
\label{app:autobounds}

Autobounds \citep{duarte2024automated} encodes the canonical model as a
polynomial program and optimizes it by spatial branch-and-bound. \Cref{thm:lp}
decides, before any solver is called, when that program is a
point or a linear program. This subsection runs both packages on the examples of
\cref{sec:illustrations} that are small enough for the general solver to
finish in seconds. The intervals agree. The difference is the time, and it
tracks the classification.

\begin{table}[ht]
\centering
\caption{One call of each package on the same table, graph, and query.
Times are wall-clock on a laptop. The Autobounds call for the mediated graph
was stopped at a limit of $25\,\mathrm{s}$; its primal and dual endpoints
agree, so the returned interval is certified. On the front door the two
Autobounds endpoints differ by about $0.002$ and both round to $0.65$.}
\label{tab:autobounds}
\small
\begin{tabular}{@{}lllr@{}}
\toprule
Example & Interval & \pid{} & Autobounds \\
\midrule
Bow, $P(Y_1{=}1)$ on~\eqref{eq:re1} & $[0.350,0.850]$ & $<1\,\mathrm{ms}$ & $0.6\,\mathrm{s}$ \\
Bow, monotone response & $[0.550,0.850]$ & $1\,\mathrm{ms}$ & $0.5\,\mathrm{s}$ \\
Markovian bow, $P(Y_1{=}1)$ & $0.700$ & $1\,\mathrm{ms}$ & $0.5\,\mathrm{s}$ \\
Instrument, seed $2$ & $[0.283,0.615]$ & $14\,\mathrm{ms}$ & $0.5\,\mathrm{s}$ \\
Front door & $0.650$ & $1\,\mathrm{ms}$ & $5.5\,\mathrm{s}$ \\
Mediated chain on~\eqref{eq:re3} & $[0.470,0.570]$ & $1\,\mathrm{ms}$ & $25\,\mathrm{s}$ \\
\bottomrule
\end{tabular}
\end{table}

On the bow, with or without monotone response, and on the instrument,
\cref{thm:lp} leaves a single simplex: one $c$-component of $\lambda=1$, or,
for the instrument, the unread component $\{Z\}$ dropped by the directed cut.
The sharp interval is then a linear program, and on the binary bow and the
binary instrument it is a closed form. Autobounds computes the same interval
by branch-and-bound over the canonical parameterization. The answers match to
three decimals, including the ATE rows omitted from \cref{tab:autobounds}
($[-0.350,0.650]$ under the order alone and $[0.000,0.650]$ under monotone
response), and the linear solver finishes in milliseconds rather than about
half a second.

The Markovian bow and the front door are the case $K=0$ of \cref{thm:lp}.
The identification algorithm of \citet{shpitser2006identification} returns a
functional of $P$, and \pid{} evaluates it. Autobounds does not
take that branch: it builds and searches a program whose optimum happens to
be a point. That search is cheap for the two-node bow and takes
$5.5\,\mathrm{s}$ for the front door.

The mediated chain is the instance where the theorem changes the program
rather than only the solver. The graph is $X\to W\to Y$ with one latent on
$\{X,W\}$, the table is~\eqref{eq:re3}, and the query is $P(Y_1{=}1)$. In the
latent laws the outcome contributes a second factor, $P(U_Y)$. The event uses
$Y$ in one world, so \cref{lem:middle} pins $Y$ to the kernel $P(Y\mid W)$
and \cref{thm:lp} leaves an LP in $P(U)$ (\cref{ex:mediated}).
\pid{} returns $[0.470,0.570]$ in about $1\,\mathrm{ms}$.
Autobounds optimizes the unreduced program and returns the same interval at
the $25\,\mathrm{s}$ limit, with the primal--dual gap closed. This is the
same graph that lies outside the symbolic class of
\citet{symbolicbounds2023general} (\cref{sec:sachs}).

The same comparison shows a difference in how assumptions are stated. Monotone
response is one named constraint, intersected with the causal order:
\begin{lstlisting}
a = Assumptions.order("X", "Y") & Assumptions.mtr("Y", "X")
bound(V, P, Query.interventional(outcome="Y", do={"X": 1}, value=1), a)
\end{lstlisting}
The Autobounds call for the same row builds the confounded bow by hand and
sets one response event to probability zero:
\begin{lstlisting}
dag = DAG("X -> Y, U -> X, U -> Y", unob="U")
pro = causalProblem(dag)
pro.add_assumption(pro.p("Y(X=1)=0&Y(X=0)=1"), "==", 0)
pro.set_estimand(pro.p("Y(X=1)=1"))
pro.load_data(summary=table)
pro.solve()
\end{lstlisting}
Further assumptions in \pid{} are further names on the same
object. \Cref{tab:catalog} is the list: an instrument, an anchored sensitivity
budget, monotone selection, a support bound. The operator \texttt{\&}
intersects them, and \cref{thm:lp} is re-checked on the intersection, so a
constraint that preserves linearity stays on the linear-program path. In
Autobounds each extra assumption is a polynomial constraint the user writes,
and the solver does not change when that constraint happens to be linear.

The comparison is confined to instances \cref{thm:lp} certifies. On a graph
the criterion rejects, such as the latent triangle in \cref{tab:audit},
\pid{} returns an outer interval with \texttt{tight=False}.
Autobounds' branch-and-bound is the general method there, and we did not run
it.

\clearpage
